% !TEX root = ../main.tex

\chapter{全文总结与展望}

\section{全文总结}

本文围绕金融大规模数据异常检测系统的实现与优化展开研究。面对传统异常检测系统在大规模场景下暴露出的异常评分复杂度高、候选召回吞吐不足、主机---设备边界交换频繁以及多 GPU 扩展效率受限等问题，论文聚焦系统组织方式，沿用已有异常检测理论模型。围绕这一边界，全文尝试将候选召回、异常评分、数据路径优化与多 GPU 扩展放入同一框架中考察，并据此验证其在金融近实时异常检测场景中的工程可行性。

在系统架构层面，论文以 PyTOD 为异常评分模块、以 FlashANNS 为候选召回前端，构建了 FlashTOD。这样组织的目的，是把候选召回和异常评分纳入同一条执行链路，使系统能够先在大规模样本中缩小评分范围，再在候选集合上执行张量化异常评分，从而缓解全量评分带来的时间与空间压力。

围绕这一执行链路，论文进一步完成了数据准备与任务建模工作。针对金融交易数据、账户行为数据和图金融数据的差异，本文对输入形式、向量化表示、候选召回目标和异常评分目标进行了统一建模，并形成了按规模逐层推进的数据组织方案。借助这一方案，全文把真实有效性验证、桥接实验、单卡性能口径和多 GPU 压力场景纳入同一实验框架。

在系统实现层面，单 GPU 场景的工作重点放在数据路径压缩上。论文围绕 GPU 主导的端到端执行链路，对查询特征组织、候选结果传递和评分中间张量的驻留方式进行了重新安排，并通过页锁定内存、RAPIDS 后端重构和 FAISS-GPU 前期接口验证降低阶段边界上的重复搬运与同步等待；完整路径中的主要候选召回前端则以 FlashANNS 为主。在这一基础上，多 GPU 部分继续沿用吞吐优先的思路，选择数据并行与索引复制作为主扩展方案，使每张 GPU 尽可能独立完成本地候选召回与异常评分；当索引复制不可行时，再以设备侧局部结果组织、NCCL 通信和设备侧最终归并作为兜底路径，并结合拓扑感知的 I/O 通道绑定与准入控制减轻共享路径冲突。

本文前期利用 FAISS-GPU 较完善的 ANN API 完成候选召回接口和 PyTOD 评分衔接验证，并在完整 FlashTOD 路径中以 FlashANNS 作为主要候选召回前端。该处理使 FAISS-GPU 的作用限定在工程验证和接口参照层面，避免将其与 FlashANNS 在大规模外存 ANNS 路径中的职责混同。

综合实验部分则把前述系统设计放回统一验证框架中。在本文设定的数据规模、硬件平台和候选预算条件下，结果显示 FlashTOD 在真实数据口径和 1M 合成桥接口径上能够保持合理的检测效果；在更大规模的高保真合成数据上，单卡优化的主要收益来自候选召回到异常评分之间的数据路径压缩，而多 GPU 数据并行与索引复制方案表现出一定的吞吐扩展能力。随着规模继续增大，I/O 路径与拓扑关系对扩展上限的影响也变得更加明显。

本文的比较证据由三个层次构成：第二章分析已有评分、检索与数据路径技术的能力边界；第六章在统一 PyTOD 评分主干下比较全量评分参考、FAISS-GPU + PyTOD 工程组合基线和 FlashTOD；单与多 GPU 路径对照用于解释性能收益的来源。这些结果能回答本文设计与代表性比较方案之间的差异，但不支持对所有 ANNS 前端或异常检测系统声称全面优势。

综合来看，本文的主要研究结论与技术贡献可以归纳为三点：
\begin{enumerate}
  \item FlashTOD 将高吞吐候选召回与异常评分组织到同一执行链路中，使系统从全量高复杂度评分转向候选缩减后的精排评分路径，并在本文实验口径下取得更稳妥的效果与效率平衡。
  \item 单 GPU 场景下，系统性能的主要限制集中在数据路径与阶段同步成本；通过 GPU 主导执行链路、RAPIDS 后端重构、FAISS-GPU 初步验证经验和 FlashANNS 候选召回前端接入，可以更稳定地改善端到端吞吐与延迟表现。
  \item 多 GPU 场景中，数据并行与索引复制可作为吞吐优先条件下的主扩展方案；当索引无法复制时，GPU 侧归并相比 CPU 聚合更有利于减少主机侧路径开销，而扩展收益能否真正保留下来，还取决于存储拓扑与 I/O 路径组织方式。
\end{enumerate}

\section{未来工作展望}

本文的工作主要集中在单机环境下金融大规模异常检测系统的实现与优化。沿着这一主线继续向前推进，仍有若干方向值得进一步深入。这些方向来自本文已经触及但尚未展开的研究边界。

当前工作已从异常评分、候选召回和系统执行路径三个层面完成代表性比较：PyTOD 全量评分提供效果参考，FAISS-GPU + PyTOD 工程组合基线控制简单 ANN 前端变量，FlashTOD 内部路径消融解释数据搬运、候选组织和多 GPU 归并开销。受实验周期和系统范围限制，本文没有穷尽所有 ANNS 前端的替换实验。后续可在统一数据、候选预算、查询集和 PyTOD 评分后端下，扩展比较 HNSW、PQ/IVF-PQ、DiskANN 与 SPANN 等候选检索路径\citep{malkov2020hnsw,jegou2011pq,subramanya2019diskann,chen2021spann}，并参照 ANN-Benchmarks 的可复现评价思路统一报告召回率、候选预算、QPS、尾延迟及资源占用\citep{aumuller2020annbenchmarks}。本文尚未完成这些替换实验，现有结果也不支持对不同 ANNS 前端作普遍优劣判断。

\subsection{面向更真实金融场景的数据验证}

本文已经在真实或半真实数据上验证了方法可用性，也借助高保真合成数据完成了大规模系统实验，但更接近生产环境的金融业务数据仍然有限。在严格遵守隐私与合规要求的前提下，后续若能接入更复杂分布、更强业务噪声和更明显概念漂移条件下的数据，将有助于进一步检验 FlashTOD 在真实业务约束下的稳定性，也能更准确地评估系统结论向生产场景外推的边界。

\subsection{面向多机场景的分布式扩展}

本文重点解决的是单机多 GPU 条件下的系统路径组织，多机场景中的跨节点通信、远程结果归并和分布式索引管理尚未进入论文范围。如果继续沿用控制面与数据平面分离的思路，把它扩展到多节点环境，就有机会进一步回答更大规模集群中的查询分发、跨节点索引一致性以及层次化归并问题，这将决定本文方案能否从单机扩展走向更大规模部署。

\subsection{面向更深层存储路径的优化}

本文已经讨论了外存驻留路径、拓扑感知 I/O 控制和 GDS 相关约束，但对更深层的存储访问优化仍然展开得不够。后续可以在现有单机链路之外评估 GPUDirect RDMA，使远端设备或存储数据在满足拓扑和驱动约束时减少主机内存中转\citep{nvidia2026gpudirectrdma}；也可参考 BaM 对 GPU 发起访问、请求批处理和 NVMe 路径的组织方式\citep{qureshi2023bam}，进一步分析查询调度与存储访问的联动。当前实现没有部署 RDMA，也没有集成 BaM。这些技术能否在 FlashTOD 中保留收益，仍需结合网络、SSD、PCIe/NUMA 拓扑和尾延迟进行独立验证。

\subsection{面向图异常检测与时序异常检测的统一扩展}

本文已将图金融数据的节点属性和邻域统计压缩为固定维度向量，但系统并未直接执行图神经网络或时序图模型。后续可参考比特币反洗钱图建模和 CARE-GNN 对交易关系、异质邻居及噪声连接的处理\citep{weber2019bitcoinaml,dou2020caregnn}，研究图表示如何进入候选召回与异常评分链路；对于随时间变化的实体关系，还可借鉴图偏差网络对多变量时间序列结构的建模方式\citep{deng2021gdn}。这类扩展会新增训练、图采样、在线特征更新和模型服务成本，不能直接沿用本文表格向量实验的性能结论。

\subsection{面向国产 GPU 平台的验证与部署}

本文当前实验与实现主要建立在现有 NVIDIA GPU 软件栈之上，而国产 GPU 平台在算子支持、向量检索实现、设备间通信机制和存储接口适配上往往存在不同约束。围绕这些差异进一步验证和移植本文方案，不仅能够检验系统设计中哪些部分具有较强的平台无关性，也有助于探索等价或近似等价的工程实现路径，从而提升方案在不同计算平台上的实际部署价值。
